\documentclass[11pt]{article}
\usepackage[utf8]{inputenc}
\usepackage{amsmath,amssymb,amsthm,pifont}
\usepackage{geometry}
\geometry{a4paper, margin=1in}

\newtheorem{proposition}{Proposition}
\newtheorem{theorem}{Theorem}
\newtheorem{lemma}{Lemma}
\newtheorem{observation}{Observation}
\theoremstyle{definition}
\newtheorem{definition}{Definition}
\theoremstyle{remark}
\newtheorem{remark}{Remark}

\title{Note 3: The PIVP Compiler Problem}
\author{Xiang Huang \and Zinan Huang}
\date{April 9, 2026}

\begin{document}
\maketitle

\section{The Problem}

Computing $\zeta(3) \in R_{\mathrm{RTCRN}}$ requires a polynomial PIVP with rational initial conditions whose output converges to $\zeta(3)$. The known ODE system for $\zeta(3)$ (via the Beukers--Apéry integral $\int_0^\infty t^2/(e^t-1)\,dt$) requires transcendental initial conditions involving $e$ and $1/(e-1)$.

\begin{definition}[PIVP Compiler]
A \textbf{PIVP compiler} is a systematic method that transforms a sequential composition of PIVPs:
\begin{enumerate}
  \item Block A: a PIVP with rational ICs computing $\alpha = \lim_{t\to\infty} u(t)$ (e.g., $\alpha = e$),
  \item Block B: a PIVP with IC depending on $\alpha$ computing $\beta = \lim_{t\to\infty} z(t)$ (e.g., $\beta = \zeta(3)$),
\end{enumerate}
into a \textbf{single} polynomial PIVP with rational ICs whose output converges to $\beta$.
\end{definition}

\section{Approaches Explored}

\subsection{Gate Chain (Time-Scale Separation)}

\textbf{Idea.} Run Block A freely. Gate Block B's evolution with a chain of logistic variables $g_k' = g_{k-1}^2(1-g_k)$. Block B only starts after Block A has approximately converged.

\textbf{Properties.}
\begin{itemize}
  \item Polynomial RHS, rational ICs, bounded trajectories. \checkmark
  \item Gate chain of depth $k$ reduces transient error geometrically. \checkmark
  \item Error is nonzero for any \emph{fixed} $k$ --- only approximate. \ding{55}
\end{itemize}

\textbf{Why it fails.} Block B computes an integral $I = \int_0^1 f(x, u(\tau(x)))\,dx$, where $u(\tau(x))$ is Block A's output at the time Block B visits $x$. Since $u \neq \alpha$ during early integration, $I \neq \int_0^1 f(x,\alpha)\,dx$. The error is
\[
  \varepsilon = \int_0^1 \bigl[u(\tau(s)) - \alpha\bigr] \cdot \frac{\partial f}{\partial u} \cdot \frac{1}{1+s}\,ds,
\]
which is bounded but nonzero for any finite gate depth.

\begin{observation}[Numerical]
For the test problem $e \cdot \ln 2$: gate depth 1 gives 13.5\% error, depth 7 gives 0.01\% error. Error $\approx C \cdot 3^{-k}$ for depth $k$.
\end{observation}

\subsection{Direct Replacement}

\textbf{Idea.} In the Apéry $\zeta(3)$ system (12 variables), replace the static constants $e, C=1/(e-1), D=e/(e-1)$ with dynamic estimators $E_1(t) \to e$, $E_2(t) \to 1/(e-1)$, $E_3(t) \to e/(e-1)$.

\textbf{Result.} Fails. The target variable $s_3(t)$ does not converge to $\zeta(3)$. ``The time-varying nature of the $E_k$ functions interacts with the main system in a non-trivial way'' (replacement.tex, DNA32 project).

\textbf{Why it fails.} Same mechanism as the gate chain: the main computation ``sees'' the wrong parameter values during its transient phase, and the accumulated error persists.

\subsection{Explosion as Infinity}

\textbf{Idea.} Use a finite-time blow-up $v' = v^2$, $v(0) = 1$ (explosion at $t=1$) to create ``infinite effective time'' for Block A. Block A converges to $\alpha$ by $t = 1^-$. Then Block B runs in $t \in (1,2)$, using a second explosion at $t=2$.

\textbf{Properties.}
\begin{itemize}
  \item $v = 1/(1-t)$, effective time $s_1 = -\ln(1-t) \to \infty$ as $t \to 1^-$.
  \item In $s_1$-time, $u$ converges to $e$. \checkmark
\end{itemize}

\textbf{Obstruction 1: Freezing.}
After DNA32 resolution ($w = 1/v = 1-t$), the variable $u(t) = e^t$ passes through $e$ at $t=1$ but does not freeze. Freezing requires a polynomial $f$ with $f(e) = 0$; since $e$ is transcendental, no such polynomial exists.

\textbf{Obstruction 2: Stability Reversal.}
For the logistic test $u' = u(1-u)/w$ with $w = 1-t$: the fixed point $u=1$ is stable for $w > 0$ (before explosion) but \emph{unstable} for $w < 0$ (after). The sign change in $w$ reverses the eigenvalue. Numerically: $u(2) \approx 10^6$.

\textbf{Obstruction 3: Snapshot Problem.}
In a concurrent ODE, Block B reads $u(t)$ (current value), not $u(1)$ (snapshot). Even with the explosion giving infinite effective time at $t=1$, Block B's integration variable sweeps $0 \to 1$ \emph{during} the explosion, and $u$ changes during this sweep.

\section{Two Types of Limits}

\begin{definition}
A PIVP computation is of \textbf{fixed-point type} if $\lim_{t\to\infty} y(t)$ converges to a globally attracting fixed point of the vector field (e.g., $u' = u(1-s)$, $s' = 1-s$ gives $u \to e$).

A PIVP computation is of \textbf{integral type} if $\lim_{t\to\infty} y(t) = \int_0^1 f(x)\,dx$, computed by accumulation along a trajectory (e.g., $I' = f(x)(1-x)$, $x' = 1-x$).
\end{definition}

\begin{observation}
Fixed-point computations are robust to parameter errors: the attractor ``washes out'' transient errors. Integral computations accumulate errors from the entire trajectory. This distinction is critical for PIVP compilation.
\end{observation}

\section{The Lyapunov Adaptive Control Approach}

\textbf{Proposed framework} (gemini\_draft.tex, DNA32 project):

\begin{enumerate}
  \item \textbf{GPAC Hierarchy (GPAC-H).} Floor 0: rational ICs. Floor $n$: uses Floor $n{-}1$ constants.
  \item \textbf{Certainty Equivalence.} Replace unknown constants with dynamic estimators.
  \item \textbf{Lyapunov Design.} Find $V$ (function of all state variables and estimation errors) with $V' \le 0$.
  \item \textbf{Dynamic Adaptation Gains.} Adaptation rates are polynomial state variables.
\end{enumerate}

\textbf{Key property.} If $V' \le 0$ can be achieved with polynomial dynamics and rational ICs, the compiled system converges to the correct value \emph{regardless of transient estimator errors}. This would unify fixed-point and integral type computations.

\textbf{Open question.} Can a suitable Lyapunov function be constructed for the 12-variable Apéry system?

\section{Summary}

\begin{center}
\renewcommand{\arraystretch}{1.3}
\begin{tabular}{lccc}
\hline
\textbf{Approach} & \textbf{Poly RHS?} & \textbf{Rational ICs?} & \textbf{Zero error?} \\
\hline
Gate chain (depth $k$) & Yes & Yes & No ($\sim 3^{-k}$) \\
Direct replacement & Yes & Yes & No (diverges) \\
Explosion + resolution & Partially & Yes & No (unfreezing) \\
Lyapunov adaptive & Yes (if exists) & Yes & Yes (if $V' \le 0$) \\
\hline
\end{tabular}
\end{center}

The PIVP compiler problem reduces to: \textbf{construct a Lyapunov function $V$ for the coupled estimator-plus-main system, with $V' \le 0$ achieved by polynomial dynamics.} Alternatively, find a fundamentally different compilation scheme (e.g., via explosions) that achieves zero error.

\end{document}
